% الجزء: sets-functions-relations
% الفصل: relations
% القسم: equivalence-relations

\documentclass[../../../include/open-logic-section]{subfiles}

\begin{document}

\olfileid[ar]{sfr}{rel}{eqv}

\olsection{علاقات التكافؤ}

تجتمع في علاقة الهوية على المجموعة ثلاث خصائص: الانعكاسية والتناظر
والتعدية. والعلاقات~${\OLId{latin-looped}{R}}$ الجامعة لهذه الثلاث كثيرة الورود.

\begin{defn}[علاقة التكافؤ]
\emph{علاقة التكافؤ} علاقة ${\OLId{latin-looped}{R}} \subseteq {\OLId{latin-looped}{A}}^{\OLMathNumeral{2}}$ تجمع الانعكاسية والتناظر
والتعدية. ويقال إن !!^{element}s ${\OLId{latin-ordinary}{x}}$ و${\OLId{latin-ordinary}{y}}$ من~${\OLId{latin-looped}{A}}$
\emph{متكافئان وفق ${\OLId{latin-looped}{R}}$} متى كان~$Rxy$.
\end{defn}

ومن علاقة التكافؤ تؤخذ \emph{فئات التكافؤ}. وذلك أن العلاقة تقسم
المجال كتلًا: تتعلق الأشياء في كل كتلة بعضها ببعض، ولا يتعلق شيء
من كتلة بشيء من أخرى. وقد ينفع أن يكون الكلام على هذه الكتل
\emph{مباشرة}؛ فلهذا نعرّفها كما يأتي.

\begin{defn}\ollabel{def:equivalenceclass}
لتكن ${\OLId{latin-looped}{R}} \subseteq {\OLId{latin-looped}{A}}^{\OLMathNumeral{2}}$ علاقة تكافؤ. إذا أخذ ${\OLId{latin-ordinary}{x}} \in {\OLId{latin-looped}{A}}$ كانت
\emph{فئة تكافؤ} ${\OLId{latin-ordinary}{x}}$ في~${\OLId{latin-looped}{A}}$ هي المجموعة $\equivrep{{\OLId{latin-ordinary}{x}}}{{\OLId{latin-looped}{R}}}
= \Setabs{{\OLId{latin-ordinary}{y}} \in {\OLId{latin-looped}{A}}}{Rxy}$. وأما \emph{مجموعة قسمة} ${\OLId{latin-looped}{A}}$ بالنسبة
إلى~${\OLId{latin-looped}{R}}$ فهي $\equivclass{{\OLId{latin-looped}{A}}}{{\OLId{latin-looped}{R}}} = \Setabs{\equivrep{{\OLId{latin-ordinary}{x}}}{{\OLId{latin-looped}{R}}}}{{\OLId{latin-ordinary}{x}} \in {\OLId{latin-looped}{A}}}$؛
أي المجموعة الجامعة لفئات التكافؤ هذه.
\end{defn}

وتشهد النتيجة الآتية لصحة هذا التعريف، إذ تبين أن فئات التكافؤ
هي حقًا الكتل التي تتجزأ إليها~${\OLId{latin-looped}{A}}$:

\begin{prop}
لتكن ${\OLId{latin-looped}{R}} \subseteq {\OLId{latin-looped}{A}}^{\OLMathNumeral{2}}$ علاقة تكافؤ. حينئذ يتكافأ $Rxy$ مع
$\equivrep{{\OLId{latin-ordinary}{x}}}{{\OLId{latin-looped}{R}}} = \equivrep{{\OLId{latin-ordinary}{y}}}{{\OLId{latin-looped}{R}}}$.
\end{prop}

\begin{proof}
أما من اليسار إلى اليمين، فنفرض $Rxy$ ونأخذ ${\OLId{latin-ordinary}{z}} \in
\equivrep{{\OLId{latin-ordinary}{x}}}{{\OLId{latin-looped}{R}}}$. فبالتعريف $Rxz$، وبما أن ${\OLId{latin-looped}{R}}$ تكافؤ يلزم $Ryz$.
وتفصيل اللزوم: من $Rxy$ وتناظر~${\OLId{latin-looped}{R}}$ نحصل على $Ryx$، ثم من
$Rxz$ وتعدية~${\OLId{latin-looped}{R}}$ نحصل على~$Ryz$. إذن ${\OLId{latin-ordinary}{z}} \in \equivrep{{\OLId{latin-ordinary}{y}}}{{\OLId{latin-looped}{R}}}$،
وبعموم الاختيار $\equivrep{{\OLId{latin-ordinary}{x}}}{{\OLId{latin-looped}{R}}} \subseteq \equivrep{{\OLId{latin-ordinary}{y}}}{{\OLId{latin-looped}{R}}}$.
وبمثل ذلك $\equivrep{{\OLId{latin-ordinary}{y}}}{{\OLId{latin-looped}{R}}} \subseteq \equivrep{{\OLId{latin-ordinary}{x}}}{{\OLId{latin-looped}{R}}}$؛ فالامتدادية
تقضي بأن $\equivrep{{\OLId{latin-ordinary}{x}}}{{\OLId{latin-looped}{R}}} = \equivrep{{\OLId{latin-ordinary}{y}}}{{\OLId{latin-looped}{R}}}$.

وأما من اليمين إلى اليسار، فنفرض $\equivrep{{\OLId{latin-ordinary}{x}}}{{\OLId{latin-looped}{R}}} =
\equivrep{{\OLId{latin-ordinary}{y}}}{{\OLId{latin-looped}{R}}}$. وانعكاسية ${\OLId{latin-looped}{R}}$ تعطي $Ryy$، ومنه
${\OLId{latin-ordinary}{y}} \in \equivrep{{\OLId{latin-ordinary}{y}}}{{\OLId{latin-looped}{R}}}$. فيكون ${\OLId{latin-ordinary}{y}} \in \equivrep{{\OLId{latin-ordinary}{x}}}{{\OLId{latin-looped}{R}}}$ أيضًا،
لأن $\equivrep{{\OLId{latin-ordinary}{x}}}{{\OLId{latin-looped}{R}}} = \equivrep{{\OLId{latin-ordinary}{y}}}{{\OLId{latin-looped}{R}}}$ بالفرض؛ أي $Rxy$.
\end{proof}

\begin{ex}
ومن أمثلة التكافؤ الحساب بترديد عدد. خذ ${\OLId{latin-ordinary}{a}}$ و${\OLId{latin-ordinary}{b}}$ و${\OLId{latin-ordinary}{n}} \in \PosInt$.
نقول ${\OLId{latin-ordinary}{a}} \equiv_{\OLId{latin-ordinary}{n}} {\OLId{latin-ordinary}{b}}$ إذا وفقط إذا اتفق باقي قسمة ${\OLId{latin-ordinary}{a}}$ على~${\OLId{latin-ordinary}{n}}$
وباقي قسمة ${\OLId{latin-ordinary}{b}}$ على~${\OLId{latin-ordinary}{n}}$. وصورته الرمزية أن ${\OLId{latin-ordinary}{a}} \equiv_{\OLId{latin-ordinary}{n}} {\OLId{latin-ordinary}{b}}$ إذا وفقط
إذا وجد ${\OLId{latin-ordinary}{k}} \in \Int$ يحقق ${\OLId{latin-ordinary}{a}} - {\OLId{latin-ordinary}{b}} = kn$. فـ$\equiv_{\OLId{latin-ordinary}{n}}$ علاقة تكافؤ
لكل~${\OLId{latin-ordinary}{n}}$، وعدد الفئات المتمايزة ${\OLId{latin-ordinary}{n}}$ بالضبط، وهي الناشئة عن~$\equiv_{\OLId{latin-ordinary}{n}}$؛ أي إن
$\equivclass{\Nat}{\equiv_{\OLId{latin-ordinary}{n}}}$ عدد !!{element}sها ${\OLId{latin-ordinary}{n}}$.
فمنها الأعداد القابلة للقسمة على ${\OLId{latin-ordinary}{n}}$ بلا باق، وهي
$\equivrep{{\OLMathNumeral{0}}}{\equiv_{\OLId{latin-ordinary}{n}}}$؛ ثم ما بقي من قسمته على ${\OLId{latin-ordinary}{n}}$ العدد~${\OLMathNumeral{1}}$،
وهو $\equivrep{{\OLMathNumeral{1}}}{\equiv_{\OLId{latin-ordinary}{n}}}$؛ و\ldots؛ إلى ما بقي من قسمته على~${\OLId{latin-ordinary}{n}}$
العدد~${\OLId{latin-ordinary}{n}}-{\OLMathNumeral{1}}$، وهو~$\equivrep{{\OLId{latin-ordinary}{n}}-{\OLMathNumeral{1}}}{\equiv_{\OLId{latin-ordinary}{n}}}$.
\end{ex}

\begin{prob}
برهن على أن $\equiv_{\OLId{latin-ordinary}{n}}$ علاقة تكافؤ لكل ${\OLId{latin-ordinary}{n}} \in \PosInt$، وأن
عدد عناصر $\equivclass{\Nat}{\equiv_{\OLId{latin-ordinary}{n}}}$ هو ${\OLId{latin-ordinary}{n}}$ بالضبط.
\end{prob}

\end{document}
